\begin{problem}[2019年第36届全国中学生物理竞赛复赛][π介子衰变与μ子观测]{109}
宇宙射线中有大量高速飞行的 $\pi$ 介子，它们会衰变成 $\mu$ 子和 $\mu$ 型反中微子 $\overline{\nu}_{\mu}$
\[
\pi \rightarrow \mu + \overline{\nu}_{\mu}
\]
已知 $\pi$ 介子和 $\mu$ 子的静止质量分别为 $m_{\pi} = 139.57061\,\mathrm{MeV}/c^{2}$ 和 $m_{\mu} = 105.65837\,\mathrm{MeV}/c^{2}$，反中微子 $\overline{\nu}_{\mu}$ 的静止质量近似为零．在实验室参照系中测得 $\pi$ 介子的飞行速度大小为 $v = 0.965c$（$c = 3.00\times 10^{8}\,\mathrm{m/s}$ 为真空中的光速）.
\begin{subproblem}
求上述衰变过程产生的 $\mu$ 子的最大速率和最小速率；
\end{subproblem}
\begin{subproblem}
上述衰变产生的 $\mu$ 子从距离地面 $10000\,\mathrm{m}$ 的高空竖直向下飞行，已知静止的 $\mu$ 子的半衰期 $T_{1/2}^{(0)} = 1.523\,\mu\mathrm{s}$，求地面观察者能够观测到的 $\mu$ 子的最大概率.
\end{subproblem}
\end{problem}